\documentclass[a4paper]{article}

\def\npart{ADAPT Short Course}
\def\nterm{Transcript-based study notes}
\def\nyear{2026}
\def\nlecturer{David Z. D'Argenio}
\def\nauthor{NanFangHong}
\def\ncourse{D'Argenio ADAPT Training Notes}
\def\ncoursehead{ADAPT Training}
\RequirePackage{etex}
\makeatletter
\ifx \npart \undefined
  \def\npart{PK/PD Notes}
\fi
\ifx \nterm \undefined
  \def\nterm{Reading notes}
\fi
\ifx \nyear \undefined
  \def\nyear{\the\year}
\fi
\ifx \nlecturer \undefined
  \def\nlecturer{the source literature}
\fi
\ifx \nauthor \undefined
  \def\nauthor{NanFangHong}
\fi
\ifx \ncourse \undefined
  \def\ncourse{Pharmacology Reading Note}
\fi
\ifx \ncoursehead \undefined
  \def\ncoursehead{\ncourse}
\fi

\author{Based on \nlecturer \\\small Notes taken by \nauthor}
\date{\nterm\ \nyear}
\title{\npart\ --- \ncourse}

\usepackage{amsfonts}
\usepackage{amsmath}
\usepackage{amssymb}
\usepackage{amsthm}
\usepackage{booktabs}
\usepackage{caption}
\usepackage{enumitem}
\usepackage{fancyhdr}
\usepackage{fontspec}
\usepackage{graphicx}
\usepackage{iftex}
\usepackage{mathtools}
\usepackage{microtype}
\usepackage{siunitx}
\usepackage{tabularx}
\usepackage{tikz}
\usepackage{xcolor}
\usepackage{xurl}
\usepackage{xeCJK}

\IfFontExistsTF{Latin Modern Roman}{
  \setmainfont[Ligatures=NoCommon]{Latin Modern Roman}
}{}

\IfFontExistsTF{Noto Serif CJK SC}{
  \setCJKmainfont{Noto Serif CJK SC}
}{
  \IfFontExistsTF{Songti SC}{
    \setCJKmainfont{Songti SC}
  }{
    \IfFontExistsTF{PingFang SC}{
      \setCJKmainfont{PingFang SC}
    }{}
  }
}

\usepackage[hidelinks,unicode,breaklinks=true]{hyperref}
\hypersetup{
  pdfauthor={\nauthor},
  pdfsubject={\npart\ - \ncourse},
  pdftitle={\npart\ - \ncourse},
  pdfkeywords={PK PD pharmacology reading notes \ncourse}
}

\pagestyle{fancyplain}
\lhead{\emph{\nouppercase{\leftmark}}}
\rhead{
  \ifnum\thepage=1
  \else
    \npart\ \ncoursehead
  \fi}

\theoremstyle{definition}
\newtheorem*{aim}{Aim}
\newtheorem*{assumption}{Assumption}
\newtheorem*{defi}{Definition}
\newtheorem*{eg}{Example}
\newtheorem*{notation}{Notation}
\newtheorem*{question}{Question}
\newtheorem*{remark}{Remark}
\newtheorem*{warning}{Warning}

\theoremstyle{plain}
\newtheorem{nthm}{Theorem}[section]
\newtheorem{nprop}[nthm]{Proposition}

\renewcommand{\labelitemi}{--}
\renewcommand{\labelitemii}{$\circ$}
\renewcommand{\labelenumi}{(\roman{*})}

\let\stdsection\section
\renewcommand\section{\newpage\stdsection}

\newcommand{\dd}{\mathop{}\!\mathrm{d}}
\newcommand{\Emax}{E_{\max}}
\newcommand{\pksym}[1]{\ifmmode\text{\textit{#1}}\else\textit{#1}\fi}
\newcommand{\pklab}[1]{\mathrm{#1}}
\newcommand{\ECfifty}{\pksym{EC}_{50}}
\newcommand{\term}[1]{\emph{#1}}
\makeatother


\usepackage{fvextra}
\usepackage{float}

\graphicspath{{adapt-short-course-figures/}}

\DefineVerbatimEnvironment{code}{Verbatim}{
  fontsize=\small,
  breaklines=true,
  breakanywhere=true,
  frame=single,
  rulecolor=\color{black!25},
  xleftmargin=0.5em,
  commandchars=\\\{\}
}

\newcommand{\concept}[1]{\textbf{#1}}
\newcommand{\lecturefigure}[2]{%
  \begin{center}
    \includegraphics[width=0.94\textwidth]{#1}
    \captionof{figure}{#2}
  \end{center}
}

\begin{document}
\maketitle
\tableofcontents

\section{ADAPT 的来处与 D'Argenio 的建模哲学}

\subsection{公开资料中的 ADAPT 背景}

ADAPT 是 USC Biomedical Simulations Resource（BMSR）长期维护的
pharmacokinetic/pharmacodynamic systems analysis software。官方页面把它定位为
面向 basic and clinical research scientists 的 computational modeling platform，
用于发现、探索和应用药物的 PK/PD properties。ADAPT 5 的用户手册作者为
David Z. D'Argenio、Alan Schumitzky 与 Xiaoning Wang。

BMSR 下载页对 ADAPT 的描述很准确：它是一组 high-level programs，用于
simulation、data analysis 和 design of experiments，主要服务于
pharmacokinetic and pharmacodynamic systems。ADAPT 5 的功能也沿着这条线展开：

\begin{itemize}
  \item \concept{SIM}: individual and multisubject simulation.
  \item \concept{ID}: individual estimation, including weighted least squares,
  maximum likelihood, generalized least squares, and MAP Bayesian estimation.
  \item \concept{SAMPLE}: optimal sampling schedule design.
  \item \concept{Population analysis}: MLEM, ITS, STS, naive pooled data and related
  nonlinear mixed-effects workflows.
\end{itemize}

1979 年 D'Argenio 与 Schumitzky 的早期论文已经把问题说得很清楚：软件要能处理
linear and nonlinear models、multiple inputs、multiple outputs、differential
equation models、non-uniform repetitive dosage regimens，并用 weighted least squares
完成 parameter estimation。今天的 ADAPT 5 只是把这一条系统建模路线继续扩展到
Bayesian estimation、population analysis、PBPK 和更复杂的 PK/PD 结构。

参考：
\begin{itemize}
  \item USC BMSR ADAPT: \url{https://bmsr.usc.edu/software/adapt/}
  \item USC BMSR downloads: \url{https://bmsr.usc.edu/downloads/}
  \item ADAPT citations: \url{https://bmsr.usc.edu/software/adapt/citations/}
  \item ADAPT 5 User's Guide: \url{https://customsitesmedia.usc.edu/wp-content/uploads/sites/106/2013/02/17062320/ADAPT5-User-Guide.pdf}
  \item D'Argenio DZ, Schumitzky A. \emph{A program package for simulation and parameter estimation in pharmacokinetic systems}. \url{https://doi.org/10.1016/0010-468X(79)90025-4}
\end{itemize}

\subsection{D'Argenio 在培训中给出的软件哲学}

D'Argenio 说 ADAPT 的开发原则是：尽量不给研究者的模型想象力设边界。这句话要认真听。
早期 ADAPT 诞生时，PK 仍是主场，PD 建模还在早期，indirect response model、TMDD、
systems pharmacology 这些今天看起来理所当然的模型类别尚未成熟。如果软件只是把当时
已有模型做成菜单，它很快会过时。

\lecturefigure{adapt-background.jpg}{ADAPT 背景与 BMSR 传统。}

因此 ADAPT 不是一个 ``choose a model from a library'' 的工具。它更像一个
dynamic systems laboratory：用户要自己写方程、自己组织 data file、自己维护 parameter
file。这个负担换来一种自由：新的 PK/PD 机制出现时，只要它能写成 state equation 和
output equation，就能进入同一套 simulation/estimation/design workflow。

\section{两天课程的主线}

\subsection{课程目标}

Day 1 的开场很重要。D'Argenio 明确说，虽然这是 ADAPT short course，但真正要学的不是
某个菜单在哪里，而是 general modeling concepts：dynamic systems、simulation、
estimation、statistical error model、debugging、population uncertainty。这些概念会
超越任何一个具体软件。

\lecturefigure{day1-agenda.jpg}{ADAPT short course 的开场与课程范围。}

Day 1 的路线是：

\begin{enumerate}
  \item ADAPT 的 systems modeling framework.
  \item Direct/indirect response simulation.
  \item TMDD simulation.
  \item Least squares, maximum likelihood, Bayesian estimation.
  \item Indirect response model 的 individual estimation.
  \item Debugging and advanced coding.
\end{enumerate}

Day 2 的路线是：

\begin{enumerate}
  \item Simple PBPK/circulatory model simulation.
  \item Minimal PBPK with naive pooled data estimation.
  \item ACI/background batch mode.
  \item Population simulation.
  \item 第二份 transcript 在 noon break 前结束；D'Argenio 说明午后会进入
  hierarchical/population modeling/nonlinear mixed effects modeling，但这段 Day 2 午后
  录音/transcript 不在当前材料中。第一份 transcript 已经覆盖 Day 1 下午的 estimation。
\end{enumerate}

\subsection{一句话总纲}

整套课可以压成一句话：

\begin{quote}
先把 pharmacology 问题写成 state equation + output equation + measurement model +
parameter model；然后根据已知和未知，选择 simulation、deconvolution、identification、
Bayesian estimation、population simulation 或 population modeling。
\end{quote}

这句话也是以后读 ADAPT、NONMEM、Monolix、SimBiology、Stan、mrgsolve 时的通用底层语言。

\section{Systems framework: input, system, output}

\subsection{三个基本问题}

D'Argenio 从工程系统讲起：飞机是 system，pilot controls 是 inputs，trajectory 是 output。
药理学也一样：dose 是 input；plasma concentration、effect-site concentration、biomarker、
clinical endpoint 是 output；中间的 absorption、distribution、metabolism、binding、
turnover、signal transduction、disease process 是 system state。

在这个框架中有三类基本问题：

\begin{enumerate}
  \item \concept{Simulation problem}: 已知 system 和 input，求 output。给定模型、参数和给药方案，预测 concentration-time profile。
  \item \concept{Deconvolution problem}: 已知 system 和 output，反推 input。典型例子是 IV 数据已刻画 PK system，口服后测得浓度曲线，想反推 absorption/input profile。
  \item \concept{Identification problem}: 已知 input 和 output，反推 system structure and parameters。在 pharmacometrics 里这就是 estimation and model selection。
\end{enumerate}

Deconvolution 在这门课里不是重点，但它很值得放在同一张地图上。它提醒我们：
``输入'' 也可以成为未知量。比如 IV bolus 数据已经给出 disposition system，口服或皮下给药后
只观察到 plasma concentration，就可以把 absorption profile 看成 hidden input。这个问题和
ordinary parameter estimation 的区别是：未知的不是几个 scalar parameters，而是一整条
time-varying input function。

这里的 identification 比 parameter estimation 更大。真实建模不是只在一个固定结构里找
参数，而是在同时问：几个 compartments？是否需要 effect compartment？是否需要 receptor
binding？是否有 saturation？是否有 delay？是否需要 covariate 或 population hierarchy？

\subsection{ADAPT 的核心模型形式}

ADAPT 接受的基本对象是 dynamic system：

\[
  \frac{\dd x(t)}{\dd t} = f\{x(t),\theta,u(t),t\}, \qquad
  y(t) = g\{x(t),\theta,u(t),t\}.
\]

其中：

\begin{itemize}
  \item \(x(t)\) 是 states: compartment amounts, tissue concentrations, receptor occupancy, signaling proteins, disease states.
  \item \(\theta\) 是 system parameters: rate constants, \pksym{CL}, \(V\), \(IC_{50}\), \(E_{\max}\), binding constants, turnover rates.
  \item \(u(t)\) 是 inputs: bolus, infusion, covariates, measured forcing functions.
  \item \(y(t)\) 是 outputs: 真正要预测或与数据比较的量。
\end{itemize}

\lecturefigure{state-output-equations.jpg}{State equation 与 output equation 是 ADAPT 的核心语法。}

最简单的一室 IV bolus 模型可以写成：

\[
  \frac{\dd X}{\dd t}=-kX,\qquad C=\frac{X}{V}.
\]

如果有 closed-form solution，也可以不写 ODE，而直接写 output equation：

\[
  C(t)=\frac{\mathrm{Dose}}{V}\exp(-kt).
\]

但 ADAPT 的边界也要记住。它不能直接处理 differential delay equations，也不处理真正的
PDE/spatial distribution model。如果 \(\dot{x}(t)\) 依赖 \(x(t-\tau)\)，或者状态是
\(x(t,r)\)，那就超出这套 ODE state-space framework。

\section{ADAPT 文件与模型层次}

\subsection{ODE 之外还有 measurement model}

初学者容易以为写完 ODE 就写完了模型。D'Argenio 特别提醒：estimation 时至少还有
measurement model：

\[
  z_{ij}=y_i(t_j;\theta)+\epsilon_{ij}.
\]

\(z\) 是实际观测，\(y\) 是模型预测，\(\epsilon\) 是 additive output error。这个 error
不只是 assay error，也可能包括 dose time error、sample time error、sample handling error、
model misspecification。你越了解 error structure，estimation 越可靠。

\subsection{参数本身也可以有模型}

在 individual simulation 中，parameters 可以只是固定常数。但在 population 情境里，
parameters 是 random variables：

\[
  \theta_i \sim \mathcal{D}(\mu,\Omega).
\]

这里的 \(\mu\) 是 population mean，\(\Omega\) 是 covariance matrix。Bayesian individual
estimation 会把这个 distribution 当作 prior；population simulation 会从它抽样生成 virtual
subjects；population modeling 则试图从数据中反推 \(\mu\)、\(\Omega\) 和 covariate effects。

\subsection{Model file, data file, parameter file}

ADAPT 工作流围绕三类文件：

\begin{itemize}
  \item \concept{Model file}: differential equations, output equations, initial conditions, variance model, secondary parameters.
  \item \concept{Data file}: dose/input history, observation times, measurements, missing values, multiple subjects or experiments.
  \item \concept{Parameter file}: simulation values, estimation initial guesses, prior mean/covariance, covariate relationships.
\end{itemize}

ADAPT 内部只认识 \(X_1,X_2,\ldots\)、\(P_1,P_2,\ldots\)、\(XP_1,XP_2,\ldots\)、\(T\)。
可读性来自显式映射：

\begin{code}
CL = P1
V  = P2
CENT = X1
CP = CENT / V

DX(1) = RIN(1) - (CL/V)*CENT
Y(1)  = CP
\end{code}

这不是装饰。大型 PBPK 或 QSP 模型里，如果不做这样的符号映射，模型很快变成不可审计的
下标泥潭。

\section{Simulation case studies}

\subsection{Indirect response simulation}

第一类例子是二室 PK 加 indirect response PD。PK 可以用 clearance parameterization：

\[
\begin{aligned}
  \frac{\dd X_1}{\dd t}
  &= R_{in}(t)-\frac{\pksym{CL}}{V_c}X_1
     -Q\left(\frac{X_1}{V_c}-\frac{X_2}{V_p}\right),\\
  \frac{\dd X_2}{\dd t}
  &= Q\left(\frac{X_1}{V_c}-\frac{X_2}{V_p}\right).
\end{aligned}
\]

PD response 是一个 endogenous process。没有药时：

\[
  \frac{\dd R}{\dd t}=k_{in}-k_{out}R.
\]

若药物抑制 production，则可写成：

\[
  \frac{\dd R}{\dd t}
  = k_{in}\left(1-\frac{I_{\max}C}{IC_{50}+C}\right)-k_{out}R.
\]

\lecturefigure{indirect-response.jpg}{二室 PK 与 indirect response PD 的 simulation case study。}

关键不是背公式，而是理解 baseline consistency。若 \(R_0\) 是 baseline response，并且无药时系统在 steady state，则：

\[
  k_{in}=k_{out}R_0.
\]

如果 initial condition 与 turnover parameters 不匹配，模型会在没有 drug effect 的情况下自己漂移。
后面任何药效解释都会被污染。

\subsection{TMDD simulation}

TMDD 的核心思想是：drug binding to target/receptor 不只是 PD 事件，也会改变 PK。Full TMDD
model 至少包含 free drug、free receptor、drug-receptor complex、receptor synthesis、
receptor degradation、binding/unbinding、complex internalization、ordinary elimination。

\lecturefigure{tmdd.jpg}{Full TMDD model 的 simulation 入口。}

D'Argenio 用 simulation 训练学生读模型：

\begin{itemize}
  \item 如果 \(k_{on}=0\)，没有 binding，drug curve 应接近普通 compartmental elimination。
  \item 曲线偏离简单指数下降的部分，来自 binding、receptor turnover 或 internalization。
  \item 改变 elimination rate 主要影响整体下降。
  \item 改变 internalization rate 往往影响 terminal phase。
  \item 只看 total drug 可能掩盖 free drug 的 TMDD 特征。
\end{itemize}

因此复杂模型不应该先拿去 estimation。先 simulation，用 controlled parameter changes 看清楚
每个参数在曲线上的手指纹。

\section{Estimation: WLS, ML, Bayesian}

\subsection{Estimation 是 inverse problem}

D'Argenio 把 estimation 称为 inverse problem：我们有 data，要反推 ``state of nature''，
也就是 system properties 或 parameters。

Weighted least squares 的目标函数是：

\[
  \hat{\theta}_{WLS}
  =
  \arg\min_{\theta}
  \sum_j w_j\{z_j-y(t_j;\theta)\}^2.
\]

若 \(w_j=1\)，就是 ordinary least squares。若知道 measurement variance，应该让 uncertainty
大的 observation 权重小，例如 \(w_j\propto 1/\sigma_j^2\)。

\subsection{Maximum likelihood 的统计含义}

Maximum likelihood 不再手动选择权重，而是写出 error distribution。若：

\[
  z_j\mid\theta \sim
  \mathcal{N}\{y(t_j;\theta),\sigma_j^2(\theta,\beta)\},
\]

则 likelihood 是所有 observation density 的乘积：

\[
  L(\theta,\beta)=\prod_j p(z_j\mid\theta,\beta).
\]

估计目标是：

\[
  (\hat{\theta},\hat{\beta})
  =
  \arg\min_{\theta,\beta} -\log L(\theta,\beta).
\]

这里 \(\beta\) 是 error variance model 的参数。D'Argenio 的关键判断是：population modeling
必须建立在 likelihood 上。因为 population modeling 要同时描述 residual error 和
between-subject variability，只靠 least squares 很难形成正确的概率结构。

\subsection{Bayesian/MAP estimation}

Bayesian estimation 的不同之处在于它不只使用当前个体的数据，也使用 prior information。
典型场景是：你已经从某个 patient population 知道 clearance 的 population mean 与 variance；
现在来了一个新病人，只有一两个浓度点。只用这个人的 sparse data 会很不稳定；Bayesian
estimation 把 prior 与 likelihood 合并。

\lecturefigure{bayesian-map.jpg}{Prior, likelihood, posterior 与 MAP estimation 的图示。}

Bayes rule 的结构是：

\[
  p(\theta\mid z)\propto p(z\mid\theta)p(\theta).
\]

如果只取 posterior distribution 的最大点，就是 MAP estimate：

\[
  \hat{\theta}_{MAP}
  =
  \arg\max_{\theta} p(z\mid\theta)p(\theta).
\]

直觉上，MAP objective = data mismatch penalty + prior deviation penalty。数据足够多时，
likelihood 主导；数据稀疏或不含某些参数信息时，estimate 会回到 prior mean 附近。

\section{Individual estimation 的输出如何读}

\subsection{Indirect response ML estimation}

回到 indirect response model，D'Argenio 用 maximum likelihood 同时估计 PK 与 PD 参数，
并为不同 outputs 设置 variance model。结果文件里最重要的不是单个 estimate，而是一组诊断：

\begin{itemize}
  \item final parameter estimates;
  \item standard errors;
  \item negative log-likelihood;
  \item number of function or ODE evaluations;
  \item secondary parameters;
  \item parameter covariance/correlation matrix;
  \item residuals;
  \item fitted versus observed profiles.
\end{itemize}

参数相关性尤其重要。如果两个 estimates 高度相关，说明数据不能把它们区分开。模型可能
过度参数化，或者实验设计没有激发足够多的 dynamics。Initial guesses 也不是小事：
非线性模型的 likelihood surface 可能有多个局部区域，optimizer 会受初值影响。

\subsection{一个实用的 estimation checklist}

\begin{enumerate}
  \item 先 simulation，再 estimation。
  \item 检查 initial condition 是否与 baseline steady state 一致。
  \item 检查 error variance model 是否符合 assay/data mechanism。
  \item 检查 initial guesses 是否在合理量级。
  \item 检查 parameter scaling 是否极端，例如 \(10^{-6}\) 和 \(10^3\) 同时出现。
  \item 用多个 initial guesses 重复估计，排除局部最优。
  \item 看 correlation matrix，而不是只看 estimates。
\end{enumerate}

\section{Debugging and advanced coding}

\subsection{Always simulate first}

D'Argenio 的 debugging 哲学非常朴素：always simulate first。不要一上来就跑 individual
或 population estimation。先确认模型在 nominal parameters 下能正常 simulation。

\lecturefigure{debugging.jpg}{Debugging and advanced programming 的开场。}

ADAPT 可以把 intermediate results 写到 screen 或 run file。复杂模型的 ODE 会被求解成百上千次；
一个中间步骤出现 not a number，就可能表现成很难理解的 solver error。因此要主动检查：

\begin{itemize}
  \item square root 里的表达式是否可能因数值误差变成负数；
  \item denominator 是否可能为零；
  \item conditional branch 是否覆盖 \(T=0\)；
  \item parameter values 是否越界；
  \item dose/input 是否真的进入了预期 compartment。
\end{itemize}

\subsection{Inverse Gaussian absorption 的调试教训}

在 inverse Gaussian absorption 例子中，输入函数在 \(T=0\) 时会除以零。模型理论上没有错，
但数值实现必须对 \(T=0\) 单独处理：

\begin{code}
IF (T .EQ. 0) THEN
  INPUT = 0
ELSE
  INPUT = formula_that_is_defined_only_for_positive_time
ENDIF
\end{code}

这个例子的意义是：debugging 不是“把错误消息消掉”，而是让模型的数学定义与数值实现一致。

\subsection{把 AUC 和 time above threshold 写成 state}

ADAPT 擅长解 ODE。因此有些看起来不是 ODE 的 quantity，可以转成 state。AUC 可以写成：

\[
  \frac{\dd \pksym{AUC}(t)}{\dd t}=C(t),\qquad \pksym{AUC}(0)=0.
\]

Time above threshold 可以写成：

\[
  \frac{\dd T_{above}}{\dd t}
  =
  \begin{cases}
    1, & C(t)>C^\ast,\\
    0, & C(t)\le C^\ast.
  \end{cases}
\]

这是一种重要的建模手感：把 summary quantity 变成 state，ADAPT 就可以在同一个 simulation
框架里追踪它。

\section{PBPK and circulatory thinking}

\subsection{PBPK 更准确地说是 circulatory model}

D'Argenio 说 PBPK 当然是 physiologically based，但更准确地说是 circulatory model。它用
circulatory system 和 anatomy 指导模型结构。血中浓度曲线不是一个抽象 compartment 的曲线，
而是 blood flow 与所有 tissues/organs 交换后的结果。

最基本的 perfusion-limited organ mass balance 是：

\[
  \frac{\dd A_o}{\dd t}=Q_oC_{art}-Q_oC_{ven,o}.
\]

若 \(A_o=V_oC_o\)，并用 partition coefficient \(R_o\) 连接 tissue concentration 与 venous
effluent concentration：

\[
  C_{ven,o}=\frac{C_o}{R_o},
\]

则：

\[
  \frac{\dd C_o}{\dd t}
  =
  \frac{Q_o}{V_o}\left(C_{art}-\frac{C_o}{R_o}\right).
\]

\lecturefigure{pbpk-mass-balance.jpg}{PBPK organ mass balance 与 perfusion-limited assumption。}

这个公式说明 PBPK 的骨架：flow、volume、partition coefficient。liver 更复杂，因为它有
hepatic artery 和 portal circulation；venous blood 是各组织流出的汇总；lung、arterial blood、
venous blood 共同闭合循环。

\subsection{PBPK 在 ADAPT 里的实现}

PBPK 模型通常一个 model file 搭配多个 data/parameter files。例如 mouse、rat、dog、human
四个 species 使用同一套结构，但 body weight、hematocrit、flows、volumes、partition
coefficients、clearance、fraction bound 不同。

\lecturefigure{pbpk-species-scaling.jpg}{PBPK case study 中 species-specific flow、volume 与 parameter file 的设置。}

D'Argenio 在演示里反复切换 case-study PDF、ADAPT run window、plot window 与输出文件。
这件事本身就是 workflow 的一部分：ADAPT 的内置 plot 足够用来快速发现趋势，但真正要比较
多个 tissues、species 或 dose levels，最好把 \texttt{.plt} 或 result files 导出到 R、Excel
或其他绘图环境中。也就是说，ADAPT 更像求解和估计引擎，而不是最终制图工具。

ADAPT 默认安装有 maximum limits，例如 differential equations 数量、system parameters 数量、
inputs、dose events、observation times。PBPK/QSP 模型容易碰到这些上限。遇到奇怪 parser
或 solver 限制时，不一定是方程错了，也可能是 installation settings 太小。

\section{Naive pooled data and batch mode}

\subsection{Minimal PBPK with naive pooled data}

Minimal PBPK case study 的目标是从多个 dose levels 的数据同时估计关键参数，例如两个
reflection coefficients 和 clearance。Naive pooled data 的假设很强：不同 dose level、
不同 replicate animals 被当作同一个 pooled system 来估计同一组参数。它忽略个体差异，
所以不是 population modeling。

\lecturefigure{naive-pooled-data.jpg}{Minimal PBPK 的 naive pooled data estimation。}

但 naive pooled data 很有用，因为复杂模型的某些参数只有在多个 dose ranges 下才可辨识：

\begin{itemize}
  \item 低剂量可能只落在线性范围，无法识别 \(IC_{50}\) 或 nonlinear binding。
  \item 高剂量可能已经饱和，反而不利于识别低浓度区的 affinity。
  \item 把多个 dose levels 放在一起，可以让数据覆盖模型的多个 dynamic regimes。
\end{itemize}

\subsection{ACI and background mode}

ADAPT 可以交互式运行，也可以 background/batch mode 运行。ACI（ADAPT command input）
文件把交互过程中需要回答的问题都写进文件：

\begin{code}
run_file_name
data_file_name
parameter_file_name
bolus_compartment
estimation_or_simulation_option
which_parameters_to_estimate
which_initial_conditions_to_estimate
which_error_parameters_to_estimate
\end{code}

当模型开发进入 production stage，尤其是同一个 executable 要重复跑多个 data files 或多个
parameter files 时，batch mode 比 GUI 点击更可靠、更可复现。

\section{Population simulation}

\subsection{Forward Monte Carlo}

Population simulation 不是 estimation，而是 forward Monte Carlo。假设你已经知道 population
distribution：

\[
  \theta_i\sim\mathrm{LogNormal}(\mu,\Omega),
\]

并知道 covariates 与 PK parameters 的关系，例如：

\[
  V_i=v_sWT_i,\qquad \pksym{CL}_i=cl_sCrCL_i.
\]

ADAPT 在 \texttt{PRIOR} function 中输入 population mean、variance、covariance。这里的
\texttt{prior} 名字来自 MAP Bayesian estimation，但同一批 mean/covariance 也可用于
population simulation。

\lecturefigure{population-simulation-setup.jpg}{Population simulation 中 model variables、covariates 与 random parameters 的设置。}

运行时要选择 normal or lognormal distribution，决定哪些 parameters random，设定 virtual
subjects 数量和 random seed。固定 seed 的意义是：如果比较两个 dosing regimens，最好用同一批
virtual subjects，减少 Monte Carlo noise。

\subsection{Mean response is not response at mean parameter}

D'Argenio 最精彩的提醒是：

\[
  \mathbb{E}\{y(t;\theta)\}\ne y(t;\mathbb{E}\{\theta\}).
\]

\lecturefigure{population-nonlinearity.jpg}{Population mean response 与 mean-parameter response 的非线性差异。}

原因是 PK/PD dynamic models 对参数通常是非线性的。即使一室模型
\[
  C(t)=C_0e^{-kt}
\]
里，\(k\) 翻倍也不会让每个时间点的浓度简单减半。到 PD、receptor binding、saturation、
TMDD 时，这个非线性更强。

这就是 population simulation 的必要性：不能只把 population mean parameters 放进模型，
然后把那条曲线当作 population mean response。

\section{复习总结}

\subsection{ADAPT 视角下的 pharmacometrics 统一语法}

两天课程真正教的是一套 pharmacometrics grammar：

\begin{enumerate}
  \item \concept{Mechanistic dynamics}: states 如何交换、结合、合成、降解、清除。
  \item \concept{Observation}: 实验测到的是 output equation 加 measurement error。
  \item \concept{Parameter uncertainty}: individual parameters 在 population 中有 distribution。
  \item \concept{Computation}: model file, data file, parameter file, SIM/ID/NPD/POP application, run file, residual file, ACI file。
  \item \concept{Decision}: simulation 和 estimation 最终服务于 dose、sampling、model revision、study design。
\end{enumerate}

\subsection{五个必须带走的句子}

\begin{enumerate}
  \item ADAPT 的灵魂是 dynamic systems，不是内置模型菜单。
  \item Simulation、estimation、deconvolution 是同一个 input-system-output 框架下的三类问题。
  \item ODE model、measurement model、parameter model 是三个不同层次，不能混为一谈。
  \item Maximum likelihood 是 individual estimation 与 population modeling 的共同地基。
  \item 对 nonlinear PK/PD 模型，\(\mathbb{E}\{y(\theta)\}\) 通常不等于 \(y\{\mathbb{E}(\theta)\}\)，所以 population simulation 不是装饰，而是理解 variability 的必要工具。
\end{enumerate}

\end{document}
